\clearpage
\fancyhead[L]{\small ANZEIGEN.}
\hypertarget{native-512-lower}{}\label{native-512-lower}
\begin{center}\small Göttingische gelehrte Anzeigen. Stück 164. Seite 1639..1640. 1808 October 13.\\[3pt]\rule{.72\linewidth}{.3pt}\end{center}

Genf. In der Bibliothèque Britannique Mai 1808 finden wir unter der Aufschrift: \emph{Sur les vingt-une dernières Comètes et les nouvelles planètes; par le Prof.} \textsc{P. Picot}, einen Auszug aus einem Briefe von \textsc{Olbers}, worin einige (bei uns hinlänglich bekannte) Nachrichten von dem letzten Cometen und von den vier neuen Planeten gegeben werden. Schätzbar ist die Zusammenstellung der parabolischen Elemente von den 21 letzten, seit 1790 erschienenen, Cometen, nach den zuverlässigsten Bestimmungen; \emph{neun} davon sind von \textsc{Olbers} selbst. Zugleich gibt dieser vortreffliche Astronom Hoffnung zu einer eignen Schrift über seine bekannte Hypothese von der Entstehung der Asteroiden. --- Ferner: \emph{Demonstration de la création immédiate de la terre en état solide, et de l'impossibilité des causes physiques pour la formation de sa figure, par Mr. l'Abbé} \textsc{Sigorgne}. Weil sowohl der von \textsc{Lacaille} am Vorgebirge der guten Hoffnung, als die verschiedenen in der nördl. Halbkugel gemessenen Breitengrade eine unregelmässige Gestalt der Erde bewiesen hätten, könnte die Erde weder flüssig noch weich gewesen sein, als sie ihre Umdrehungsbewegung erhielt; sie hätte schon damals ein fester Körper, mithin auch schon vorher an den Polen abgeplattet sein müssen. Rec. hält weder die Prämissen für entschieden, noch die Schlussfolge für richtig. Der \textsc{Lacaille}'sche Breitengrad möchte wohl eben so sehr einer neuen Verbesserung bedürftig sein, als der Lappländische war, selbst der in Peru gemessene ist noch Zweifeln unterworfen. Wenn gleich indess sich nicht bezweifeln lässt, dass der Erdkörper in gewissem Grade unregelmässig ist, so kann man hieraus doch keinesweges auf die Unmöglichkeit einer Bildung aus physischen Ursachen schliessen; ein ganz regelmässiger Körper hätte die Erde nur dann werden können, wenn alle Theile der Masse, woraus sie sich bildete, vollkommen flüssig gewesen und so lange geblieben wären, bis sie sich nach hydrostatischen Gesetzen geordnet hätten.
